\begin{table}[htbp]
\centering
\small
\renewcommand{\arraystretch}{0.92}
\caption{符号说明}
\label{tab:symbols}
\begin{tabularx}{0.94\textwidth}{>{\centering\arraybackslash}p{0.24\textwidth} X}
\toprule
符号 & 含义 \\
\midrule
\(\mathcal I\) & 原始用频计划集合，包含 150 份计划。\\
\(i,j\) & 用频计划或装备的索引。\\
\(k\) & 同一计划内重复事件的序号。\\
\([f_i^-,f_i^+)\) & 计划 \(i\) 的频段半开区间，端点按离散频率单位记录。\\
\([s_i,e_i)\) & 计划 \(i\) 的首次时间半开区间，端点按离散时间单位记录。\\
\(l_i\) & 计划 \(i\) 的单次使用时长，满足 \(l_i=e_i-s_i\)。\\
\(g_i\) & 计划 \(i\) 的空闲间隔时长。\\
\(n_i\) & 计划 \(i\) 的重复使用次数。\\
\(E_i\) & 计划 \(i\) 展开后的完整重复事件集合。\\
\(C_{ij}\) & 计划 \(i\) 与计划 \(j\) 是否存在时频冲突的指示量。\\
\(\mathcal A_i\) & 计划 \(i\) 的候选动作集合，包括保留、频移、时移和撤销。\\
\(a\) & 候选动作索引。\\
\(x_{ia}\) & 计划 \(i\) 是否选择动作 \(a\) 的 0--1 决策变量。\\
\(\mathcal K_{ia}\) & 候选动作 \((i,a)\) 占用的离散时频资源格集合。\\
\(c=(t,f)\) & 离散时频资源格索引。\\
\(R\) & 撤销计划总数。\\
\(R_A,R_B\) & A 类、B 类撤销计划数。\\
\(M\) & 被调整的计划总数。\\
\(M_A,M_B\) & A 类、B 类被调整的计划数。\\
\(S_{10}\) & 已闭合优先级前缀下的归一化平移量整数表示；当前值为 775，但尚未证明其全局最小。\\
\(\delta f,\delta t\) & 频率平移量和首次时间平移量，均按离散资源单位记录。\\
\(\delta g\) & C 类计划的间隔调整量，决定 \(g_i'=8+\delta g\)。\\
\(p\) & 问题三中的一个完整 C 类候选起点。\\
\(y_p\) & 是否选择候选计划 \(p\) 的 0--1 决策变量。\\
\(\mathcal K_p\) & 候选计划 \(p\) 的完整资源格集合。\\
\(O_c\) & 固定问题二方案对资源格 \(c\) 的占用指示量。\\
\(H\) & 统一时间域的右端点，基准值为 643，用于边界敏感性分析。\\
\(S_{10}^{(4)}\) & 问题四在固定最小撤销层上的末级变化量候选指标。\\
\bottomrule
\end{tabularx}
\end{table}
